% ============================================================
% f2013_shaanxi_science_q12_views.tex
% 2013 陕西卷（理工类）第 12 题 三视图（TikZ 矢量重绘）
% 主视图：直角边 1、2 的直角三角形
% 左视图：底边 2、高 2 的等腰三角形
% 俯视图：半径 1 的半圆
% 还原为底面半径 1、高 2 的半圆锥，体积 π/3
% ============================================================
\begin{tikzpicture}[
  x=1cm, y=1cm, line width=0.5pt, font=\footnotesize, >=Stealth
]

  % ---------- 主视图 ----------
  \draw (0,0) -- (1,0) -- (0,2) -- cycle;
  \draw[|<->|] (-0.35,0) -- (-0.35,2) node[midway, left] {$2$};
  \foreach \y in {0,2} \draw (-0.50,\y) -- (-0.20,\y);
  \draw[|<->|] (0,-0.35) -- (1,-0.35) node[midway, below] {$1$};
  \foreach \x in {0,1} \draw (\x,-0.50) -- (\x,-0.20);
  \node at (0.5,-1.00) {主视图};

  % ---------- 左视图 ----------
  \begin{scope}[xshift=2.6cm]
    \draw (0,0) -- (2,0) -- (1,2) -- cycle;
    \draw[|<->|] (0,-0.35) -- (1,-0.35) node[midway, below] {$1$};
    \draw[|<->|] (1,-0.35) -- (2,-0.35) node[midway, below] {$1$};
    \foreach \x in {0,1,2} \draw (\x,-0.50) -- (\x,-0.20);
    \node at (1,-1.00) {左视图};
  \end{scope}

  % ---------- 俯视图 ----------
  \draw (0.5,-3.20) -- (0.5,-5.20);
  \draw (0.5,-3.20) arc (90:-90:1);
  \node at (1.20,-5.65) {俯视图};

\end{tikzpicture}
